
\newpage
\section{Общее описание базовых блоков методики}
\label{part:methods_intro}

Что же, пришло время для того, чтобы сформулировать и описать предлагаемую в работе методику анализа прибрежных территорий.
Ранее, в цели исследования, формулировалась необходимость разработки метода оценки риска проявлений негативных гидромеханических и геомеханических процессов, проявляющихся на прибрежных территориях вследствие их соседства с морем.

Исходя из этого, можно зафиксировать начальный и конечный шаги разрабатываемого алгоритма:

\begin{itemize}
\item исходными данными будет являться, задаваемая пользователем, область, гарантировано включающая в себя береговую линию;
\item выходными -- рассчитанные для этой береговой линии показатели, допускающие интерпретацию по степени риска.
\end{itemize}

Несмотря на свободу формулировок границ разрабатываемой методики они уже позволяют наложить на область возможных решений существенные и важные ограничения.
Во-первых, базовым объектом, для которого выполняется анализ, является полилиния, априорно предполагающая изменение показателей уязвимости и риска вдоль своей трассы.
Оставляя пока сам механизм расчета как некоторый абстрактный детерминированный алгоритм, логично определить причину вариативности получаемых показателей в изменении исходных данных.
Даже при константных внешних факторах, сама, изменяющая свое направление, трасса береговой линии, будет являться источником для получения различных вычисляемых значений, что не допускает единственного расчета, применимого ко всей линии целиком.

Исходя из этого, наиболее простым и эффективным решением будет выполнение необходимых расчетов для множества точек, принадлежащих исследуемой трассе с последующей интерполяцией полученных результатов.
Несмотря на кажущуюся простоту задачи выделения точек, от результатов этого этапа будет напрямую зависеть баланс между репрезентативностью получаемых результатов и вычислительной эффективностью методики.
Очевидно, что метод отбора наиболее показательных и характерных точек для расчета должен зависеть от масштаба и размера области интереса, условий местности, формы исследуемой линии и целей, определяемых решаемой задачи.

Последнее замечание особенно важно не столько в контексте выбора точек для расчета, сколько в формулировании главного вопроса стоящего перед методикой: \enquote{А собственно, что мы хотим в итоге посчитать и как интерпретировать полученные результаты?}
Отметим так же то, что сама форма исследуемой береговой линии может определяться как естественным побережьем, так и уже существующими (или проектируемыми) прибрежными инженерными конструкциями, что будет качественно влиять на специфичные объектам факторы риска и целевые совокупности исследуемых точек.

Для естественного берега и неурбанизированных территорий ответ на вопрос оценки риска уязвимости прибрежных зон был представлен группой методов по расчету CVI (описанных в рамках 1 этапа НИР).
В свою очередь для инженерных объектов (молы, набережные, причалы, пляжи и~т.\,д.) такой подход может выдавать искаженные и трудно интерпретируемые результаты.
Так, если мы хотим оценить уязвимость набережной, так ли важен для нас естественный уклон, который был или будет трансформирован при ее сооружении? 
Повлияют ли геологические свойства пород, слагающих прибрежные территории на бетонный пирс? 
Учтет ли CVI вдольбереговую абразию пляжного материала?
Естественное ограничение любой универсальной методики, примером которой является исходный CVI, не позволяет ответить уверенно и утвердительно как на приведенные вопросы, так и на вопросы, которые неизбежно поставит перед будущими аналитиками время.
Однако уже из них становятся видна необходимость возможности расширения методики оценки риска прибрежных территорий различными методами анализа, специфичными конкретному типу решаемому задачи.

Допуская необходимость перспективного использования различных подходов к решению задач оценки прибрежных территорий, и неоднозначность присущих им методов анализа, синтезированное решение целесообразно представить в виде последовательности состоящей из четырех основных блоков:
\begin{itemize}
\item блока формулирования задачи, в котором определяется предметная область расчета и формирются данные, необходимые для математического описания условий решаемой задачи;
\item блока получения, проверки и предварительной обработки исходных данных;
\item блока расчета, непосредственно определяющего внутреннюю логику работы модели и реализующего алгоритмы расчета показателей и критериев оценки риска;
\item блок представления информации, в котором происходит интерпретация, сохранение и визуализация полученных результатов, служащий основой для принятия управленческих и проектировочных решений.
\end{itemize}

Несмотря на определяющую и очевидную важность в разрабатываемой методике блока расчета, его четкое описание и устойчивая работа не возможны без описания логики работы окружающих его блоков.
Исходя из этого, раскорем подробнее опорные требования к функциональности каждого из них и интерфейсной связи между ними.
